\textbf{结论名称：}三角形内角平分线定理及逆定理。

\textbf{条件：}非退化 $\triangle ABC$，$D$ 严格位于线段 $BC$ 内部，所有长度为普通正长度。
\[
 \angle BAD=\angle DAC\quad\Longleftrightarrow\quad
 \frac{BD}{DC}=\frac{AB}{AC}.
\]
由 $BD+DC=BC$，得到
\[
 BD=\frac{BC\cdot AB}{AB+AC},\qquad DC=\frac{BC\cdot AC}{AB+AC}.
\]
\textbf{低门槛证明：}角平分线上的点到两邻边所在直线距离相等；
同一对三角形 $ABD,ACD$ 的面积比，按共同高算为 $BD/DC$，按相等的高算为 $AB/AC$。
逆向另作真正的内角平分线，再利用内部比例分点的唯一性。

\textbf{例：}$AB=6,AC=4,BC=5$ 时，$BD=3,DC=2$。
三边满足 $|6-4|<5<6+4$，且 $3+2=5,3/2=6/4$。
一般情形不能把角平分线当成中线或高；在它已是角平分线时，三者重合当且仅当 $AB=AC$。
